\chapter{Proofs}
\label{app:proofs}

This appendix records the proof-facing structure of the dissertation. It is not intended to
replace the theorem chapter; it is intended to make the proof obligations inspectable and to
show which claims are formal, which are empirical, and which rely on shared typed runtime
objects.

\section{Proof map}
\begin{longtable}{p{0.18\textwidth}p{0.28\textwidth}p{0.18\textwidth}p{0.24\textwidth}}
\caption{Theorem-to-proof map for the ORIUS dissertation.}\label{tab:proof-map}\\
\toprule
\textbf{Result} & \textbf{Main idea} & \textbf{Depends on} & \textbf{Artifact-facing link}\\
\midrule
\endfirsthead
\toprule
\textbf{Result} & \textbf{Main idea} & \textbf{Depends on} & \textbf{Artifact-facing link}\\
\midrule
\endhead
Theorem~\ref{thm:oasg-existence} & If observed legality is not informationally sufficient for true-state safety, then there exists a degraded-observation regime that admits an observed-safe but true-unsafe release. & A1, A4, A5 & Dual-trajectory replay and the OASG metric family.\\
Theorem~\ref{thm:safety-preservation} & If inflation is conservative and the repair or fallback operator respects the tightened set, the released action preserves the defended safety predicate. & A1, A2, A3, A5 & Repair traces, fallback logs, and certificate-valid release surfaces.\\
Theorem~\ref{thm:risk-envelope} & As reliability falls, uncertainty inflation and admissible-set tightening contract the release envelope lawfully rather than heuristically. & A1, A3, A6 & Calibration diagnostics, theorem gate figure, and runtime intervention traces.\\
Proposition~\ref{prop:no-free-safety} & Quality-ignorant release can remain observed-safe while incurring true-state failure. & A1, A4, A5 & Witness-row violation tables and fault-family replay.\\
T10/T11 sandwich & Observation-only mandatory release has a Bayes ambiguity-risk lower bound, while ORIUS has an $\alpha$-covered upper bound under correct uncertainty-set coverage. & A1--A6 plus coverage, safe-set correctness, repair membership, fallback admissibility, and adapter validity & Observation-ambiguity module, focused tests, and promotion matrix.\\
T5 & A certificate is valid for a finite horizon when the reachable tube remains inside the safe set and no invalidating event occurs. & A1--A5 plus reachable-tube containment & Temporal-theorem code, horizon artifact, and result card.\\
T8 & ORIUS weakly dominates blind persistence in true-state violations while retaining a stated fraction of useful work under the same fault trace. & A1--A7 plus paired-trace comparability & Graceful-degradation module, policy comparison artifact, and result card.\\
T9/T10 & Empty common safe core implies mandatory-release impossibility; two indistinguishable boundary states imply the TV lower bound. & Mandatory release, disjoint safe sets, two-state boundary model & Ambiguity and boundary-lower-bound modules plus result cards.\\
T11Byz/Tstale & Robust reliability aggregation is valid for $b<n/2$ trimmed channels; stale uncertainty grows under bounded drift. & Honest-channel interval bound; bounded latent drift & Robust reliability and stale-uncertainty modules plus result cards.\\
L1--L4/minimax/sensor & Runtime monotonicity lemmas, scoped finite ambiguity minimax, and sensor necessity under adapter semantics. & Narrow runtime-law assumptions and finite ambiguity class & Runtime-law, minimax, and sensor-necessity modules plus result cards.\\
\bottomrule
\end{longtable}

\section{Proof sketches}
\paragraph{OASG existence.}
The proof proceeds by separating observed-state legality from true-state legality. Once the
adapter exposes a true-state violation predicate and the replay harness exposes an observed
state that is not informationally sufficient for that predicate, one can construct a
degradation regime in which observed legality survives while true-state safety fails. The
scientific role of the theorem is not to characterize every such regime exhaustively. It is
to show that the gap is structural rather than anecdotal.

\paragraph{Safety preservation under repair.}
The second theorem uses the conservativeness of the inflated observation-consistent set. If
the repair operator only emits actions admissible over that set, or the fallback operator
replaces release when such action does not exist, then the released action remains safe for
the latent state whenever the inflation really upper-bounds that state. This is the exact
point at which fallback becomes part of the proof surface rather than an implementation
afterthought.

\paragraph{Degradation-sensitive risk envelope.}
The third theorem formalizes the intuition that reduced trust must shrink release freedom.
Monotone reliability loss implies monotone widening of the uncertainty set; that widened set
then implies monotone tightening of admissible release. The theorem does not say
intervention is free. It says intervention pressure is a lawful consequence of degraded
observation.

\paragraph{No-free-safety proposition.}
The proposition follows by contradiction against any controller that checks legality only on
observed state while ignoring observation degradation. If the degraded observation regime is
admissible yet insufficiently informative about the latent state, then there exists a
release that remains locally coherent while still violating the true safety predicate.

\paragraph{Covered observation-ambiguity sandwich.}
Let $L(a,o)=\mathbb{P}(a\notin C(X^\star)\mid O=o)$ be the conditional
true-state violation risk of releasing action $a$ after observing $o$.  Any
observation-only mandatory-release controller $\pi$ selects a single action
$\pi(o)$ on the entire ambiguity class $B(o)=\{x:h(x)=o\}$.  Therefore, for
each observation value $o$,
\[
  L(\pi(o),o)\ge \min_{a\in\mathcal{A}} L(a,o).
\]
Taking expectation over $O$ gives
\[
  \operatorname{TSVR}(\pi)
  = \mathbb{E}_{O}\!\left[L(\pi(O),O)\right]
  \ge
  \mathbb{E}_{O}\!\left[
    \min_{a\in\mathcal{A}}
    \mathbb{P}(a\notin C(X^\star)\mid O)
  \right].
\]
This proves the Bayes lower bound.  If, for some $o$, the common safe core
$K(o)=\cap_{x\in B(o)}C(x)$ is empty and the conditional law of $X^\star$ puts
positive mass on states that exclude each candidate action, then the conditional
minimum is strictly positive.  If $K(o)$ is nonempty, the lower bound correctly
allows zero conditional risk; this is why the theorem is not the false claim
that all ambiguity forces violation.

For the ORIUS upper bound, suppose the released action $u_t$ satisfies
$u_t\in C(x)$ for every $x\in\mathcal{X}_t$ and
$\mathbb{P}(X_t^\star\notin\mathcal{X}_t)\le\alpha$.  Split the violation event:
\[
  \mathbb{P}(u_t\notin C(X_t^\star))
  =
  \mathbb{P}(u_t\notin C(X_t^\star),\, X_t^\star\in\mathcal{X}_t)
  +
  \mathbb{P}(u_t\notin C(X_t^\star),\, X_t^\star\notin\mathcal{X}_t).
\]
The first term is zero by set-wise release safety.  The second term is bounded
by $\mathbb{P}(X_t^\star\notin\mathcal{X}_t)\le\alpha$.  Hence
$\operatorname{TSVR}_{\mathrm{ORIUS}}\le\alpha$, and exact coverage gives zero
one-step true-state violation probability.

\section{Interpretation}
These proofs are deliberately modest in scope. They do not prove broader future-domain
closure and they do not replace the need for replay, calibration, runtime, and governance artifacts.
Their role is to anchor the manuscript's strongest claims to explicit assumptions and typed
objects so that later empirical chapters can be read as validation of the same runtime
contract rather than as isolated case studies.
